\subsection{WEIGHTS AND SPECIAL WEIGHTS}

PROPOSITION 3. The special points (Chap. V, § 3, no. 10, Def. 1) of \(W_{a}\) are the weights of \(R^{\prime}\) .

Let \(x_0 \in E\) and let \(\alpha \in R\) . The hyperplane L parallel to Ker \(\alpha\) and passing through \(x_0\) has equation \(\langle \alpha, x \rangle = \langle \alpha, x_0 \rangle\) . To be equal to some \(L_{\beta,k}\) , it is necessary on the one hand that \(\alpha\) and \(\beta\) are proportional, and so, since R is reduced, that \(\beta = \pm \alpha\) , and on the other hand that \(\langle \alpha, x_0 \rangle\) is an integer. It follows immediately that \(x_0\) is a special point of \(W_a\) if and only if \(\langle \alpha, x_0 \rangle \in Z\) for all \(\alpha \in R\) , in other words, if and only if \(x_0 \in P(R^{\prime})\) (§ 1, no. 9).

COROLLARY. (i) If \(\overline{\omega} \in \mathrm{P}(\mathbb{R}^{\prime})\) , there exists an alcove C such that \(\overline{\omega}\) is an extremal point of \(\overline{\mathbb{C}}\) .

\begin{enumerate}
\def\labelenumi{(\roman{enumi})}
\setcounter{enumi}{1}
\tightlist
\item
  If \(\mathbf{C}\) is an alcove, \(\overline{\mathbf{C}} \cap \mathbf{Q}(\mathbb{R}^{\prime})\) reduces to a single point and this is an extremal point of \(\overline{\mathbf{C}}\) .
\end{enumerate}

This follows from Prop. 3, in view of the Cor. of Prop. 11 of Chap. V, § 3, no. 10 and Prop. 12 of Chap. V, § 3, no. 10.

PROPOSITION 4. Let \(C'\) be a chamber of \(R'\) .

\begin{enumerate}
\def\labelenumi{(\roman{enumi})}
\item
  There exists a unique alcove C contained in \(\mathbf{C}'\) such that \(0\in \overline{\mathbf{C}}\)
\item
  The union of the \(w(\overline{\mathbf{C}})\) for \(w\in \mathrm{W}\) is a neighbourhood of 0 in E.
\item
  Every wall of \(\mathbf{C}'\) is a wall of \(\mathbf{C}\) .
\end{enumerate}

This follows from Prop. 11 of Chap. V, § 3, no. 10.

Assume now that \(R\) is irreducible. Let \((\alpha_i)_{i\in I}\) be a basis of \(R\) (§1, no. 5, Def. 2), and let \((\overline{\omega}_i)_{i\in I}\) be the dual basis. The \(\overline{\omega}_i\) are the fundamental weights of \(R^{\prime}\) for the chamber \(D'\) of \(R^{\prime}\) corresponding to the basis \((\alpha_i)\) . Let

\[
\tilde {\alpha} = \sum_ {i \in \mathrm{I}} n _ {i} \alpha_ {i}
\]

be the highest root of \(\mathbf{R}\) (§ 1, no. 8), and let \(\mathbf{J}\) be the set of \(i \in \mathbf{I}\) such that \(n_i = 1\) .

PROPOSITION 5. Let \(\mathbf{C}\) be the alcove contained in \(\mathbf{C}'\) and containing 0 in its closure (Prop. 4).

\begin{enumerate}
\def\labelenumi{(\roman{enumi})}
\tightlist
\item
  C is the set of \(x \in E\) such that \(\langle\alpha_{i}, x\rangle > 0\) for all \(i \in I\) and \(\langle\tilde{\alpha}, x\rangle < 1\) .
\item
  The set \(\overline{C} \cap P(R^{\prime})\) consists of 0 and the \(\overline{\omega}_{i}\) for \(i \in J\) .
\end{enumerate}

Let D be the set of \(x \in E\) such that \(\langle\tilde{\alpha}, x\rangle < 1\) and set \(C_{1} = C' \cap D\) . Since \(0 \in \overline{C}\) , we have \(C \subseteq D\) and hence \(C \subseteq C_{1}\) . We are going to show that, for all \(\alpha \in R\) and all \(k \in Z\) , the sets C and \(C_{1}\) are on the same side of the hyperplane \(L_{\alpha,k}\) . This will prove that \(C_{1} \subseteq C\) and so will establish assertion (i). If k = 0, the whole of the chamber \(C'\) is on one side of \(L_{\alpha,0}\) , which establishes our assertion in this case. If \(k \neq 0\) , we may, by replacing \(\alpha\) by \(-\alpha\) , assume that k \textgreater{} 0. Then \(\langle\alpha, x\rangle < k\) on C, since \(0 \in \overline{C}\) . On the other hand, \(\tilde{\alpha} - \alpha\) is positive on \(C'\) (§ 1, no. 8, Prop. 25). Thus, for \(y \in C_{1}\) , we have \(\langle\alpha, y\rangle \leqslant \langle\tilde{\alpha}, y\rangle < 1 \leqslant k\) . Consequently, C and \(C_{1}\) are on the same side of \(L_{\alpha,k}\) .

Now let \(\overline{\omega} \in P(R^{\prime})\) . Then \(\overline{\omega} = \sum_{i} p_i \overline{\omega}_i\) , with \(p_i \in \mathbf{Z}\) (§ 1, no. 10), and \(\overline{\omega} \in \overline{C'}\) if and only if the integers \(p_i\) are positive. If \(\overline{\omega} \in \overline{C'}\) , then \(\overline{\omega} \in \overline{C}\) if and only if \(\langle \tilde{\alpha}, \overline{\omega} \rangle = \sum_{i} n_i p_i\) is \(\leqslant 1\) , hence (ii).

COROLLARY. The alcove C is an open simplex with vertices 0 and the \(\overline{\omega}_i / n_i\) , \(i \in I\) .

This follows from (i).

